\begin{lemma}\label{src:C.4.9}
设 $\mathcal X$ 为 $\mathrm{Art}_{\Bbbk}^{\mathrm{op}}$ 上的预叠, 群胚 $\mathcal X(\Bbbk)$ 等价于集合 $\{x_0\}$; 设 $F_{\mathcal X}:\mathrm{Art}_{\Bbbk}\to\Sets$ 为将 $A\in\mathrm{Art}_{\Bbbk}$ 映到同构类集合 $\mathcal X(A)/{\sim}$ 的协变函子. 假设 $\mathcal X$ 满足条件 $(\mathrm{RS}_2)$.
\begin{enumerate}[(1)]
\item 切空间 $T_{\mathcal X}=F_{\mathcal X}(\Bbbk[\epsilon])$ 具有自然的 $\Bbbk$-向量空间结构. 更一般地, 对每个有限维 $\Bbbk$-向量空间 $V$, 用 $\Bbbk[V]$ 表示由 $V^2=0$ 定义的 $\Bbbk$-代数 $\Bbbk\oplus V$, 则集合 $F_{\mathcal X}(\Bbbk[V])$ 具有自然的 $\Bbbk$-向量空间结构, 且有函子性双射 $F_{\mathcal X}(\Bbbk[V])=T_{\mathcal X}\otimes_{\Bbbk}V$.
\item 考虑 $\mathrm{Art}_{\Bbbk}$ 中核 $I$ 的平方为零的满射 $A\twoheadrightarrow A_0$, 以及元素 $x_0\in\mathcal X(A_0)$; 用 $\operatorname{Lift}_x(A)$ 表示 $\Spec A_0\to\Spec A$ 上的态射 $\alpha:x_0\to y$ 的集合, 其中若存在同构 $\beta:x\to x'$ 满足 $\alpha'=\beta\circ\alpha$, 则令 $(\alpha:x_0\to x)\sim(\alpha':x_0\to x')$. 在 $\operatorname{Lift}_{x_0}(A)$ 上存在 $T_{\mathcal X}\otimes I$ 的作用, 它对 $\mathcal X$ 是函子性的. 假设 $\operatorname{Lift}_{x_0}(A)$ 非空. 若 $\mathcal X$ 满足条件 $(\mathrm{RS}_1)$, 则此作用传递; 若满足条件 $(\mathrm{RS}_4)$, 则此作用自由且传递 (即 $\operatorname{Lift}_{x_0}(A)$ 为 $T_{\mathcal X}\otimes I$ 下的 torsor).
\end{enumerate}
\end{lemma}
\begin{proof}
首先注意, 若 $V$ 为有限维向量空间, 则 $\Bbbk[V]=\Bbbk[\epsilon]\times_{\Bbbk}\cdots\times_{\Bbbk}\Bbbk[\epsilon]$; 归纳地应用 $(\mathrm{RS}_2)$, 可知 $(\mathrm{RS}_2)$ 的结论在 $A_0=\Bbbk$ 且 $A=\Bbbk[V]$ 时也成立. 对 $B_0\in\mathrm{Art}_{\Bbbk}$, $(\mathrm{RS}_2)$ 的第一部分蕴含 $F_{\mathcal X}(B_0\times_{\Bbbk}\Bbbk[V])\xrightarrow{\sim}F_{\mathcal X}(B_0)\times F_{\mathcal X}(\Bbbk[V])$ 为双射. 特别地, 对每对有限维向量空间, $F_{\mathcal X}(\Bbbk[V]\times_{\Bbbk}\Bbbk[W])\xrightarrow{\sim}F_{\mathcal X}(\Bbbk[V])\times F_{\mathcal X}(\Bbbk[W])$ 双射; 换言之, 函子 $V\mapsto F_{\mathcal X}(\Bbbk[V])$ 与有限积交换. $T_{\mathcal X}=F_{\mathcal X}(\Bbbk[\epsilon])$ 的向量空间结构由
\begin{equation}\label{src:C.4.10}
F_{\mathcal X}(\Bbbk[\epsilon]\times_{\Bbbk}\Bbbk[\epsilon'])\xrightarrow{\sim}F_{\mathcal X}(\Bbbk[\epsilon])\times F_{\mathcal X}(\Bbbk[\epsilon'])
\end{equation}
的双射性推出. 事实上, 若 $\tau_1,\tau_2\in F_{\mathcal X}(\Bbbk[\epsilon])$, 则可用 \bookeqref{C.4.10} 将 $(\tau_1,\tau_2)$ 看成 $F_{\mathcal X}(\Bbbk[\epsilon]\times_{\Bbbk}\Bbbk[\epsilon'])$ 的元素, 并将 $\tau_1+\tau_2$ 定义为 $(\tau_1,\tau_2)$ 在 $F(\Bbbk[\epsilon]\times_{\Bbbk}\Bbbk[\epsilon'])\to F(\Bbbk[\epsilon])$ 下的像, 此映射由将 $(\epsilon,0)$ 和 $(0,\epsilon')$ 都映到 $\epsilon$ 的环映射 $\Bbbk[\epsilon]\times_{\Bbbk}\Bbbk[\epsilon']\to\Bbbk[\epsilon]$ 诱导. $c\in\Bbbk$ 对 $\tau\in F_{\mathcal X}(\Bbbk[\epsilon])$ 的数乘, 定义为 $\tau$ 在 $F_{\mathcal X}(\Bbbk[\epsilon])\to F_{\mathcal X}(\Bbbk[\epsilon])$ 下的像, 此映射由将 $\epsilon$ 映到 $c\epsilon$ 的映射 $\Bbbk[\epsilon]\to\Bbbk[\epsilon]$ 诱导.

同一论证表明 $V\mapsto F_{\mathcal X}(\Bbbk[V])$ 定义 $\Bbbk$-线性函子 $\mathrm{Vect}_{\Bbbk}^{\mathrm{fd}}\to\mathrm{Vect}_{\Bbbk}$. 自然映射为
\[
\begin{aligned}
F_{\mathcal X}(\Bbbk[\epsilon])\times\operatorname{Hom}_{\Bbbk}(\Bbbk[\epsilon],\Bbbk[V])&\to F_{\mathcal X}(\Bbbk[V]),\\
(\tau,\phi)&\mapsto\phi^*\tau:=F_{\mathcal X}(\phi)(\tau).
\end{aligned}
\]
对于 (2), 注意自然映射
\[
A\times_{A_0}A\to\Bbbk[I]\times_{\Bbbk}A,\qquad(a_1,a_2)\mapsto(\overline a_1+a_2-a_1,a_1)
\]
为同构. 因而有图
\[
\begin{aligned}
\mathcal X(\Bbbk[I])\times\mathcal X(A)&\longleftarrow\mathcal X(\Bbbk[I]\times_{\Bbbk}A)\\
&\simeq\mathcal X(A\times_{A_0}A)\xrightarrow{p_1^*}\mathcal X(A),
\end{aligned}
\]
其中最左函子由 $(\mathrm{RS}_2)$ 的第一部分本质满射. 为定义作用, 取 $\tau\in T_{\mathcal X}\otimes I=F_{\mathcal X}(\Bbbk[I])$ 和 $(\alpha:x_0\to x)\in\operatorname{Lift}_{x_0}(A)$. 选取 $\tau$ 的代表 $\widetilde\tau\in\mathcal X(\Bbbk[I])$. 将 $\tau\cdot(\alpha:x_0\to x)\in\operatorname{Lift}_{x_0}(A)$ 定义为 $p_1^*y$, 其中 $y\in\mathcal X(\Bbbk[I]\times_{\Bbbk}A)$ 为 $(\widetilde\tau,x)$ 的原像. 为证明良定义, 假设 $y'\in\mathcal X(\Bbbk[I]\times_{\Bbbk}A)$ 为另一原像. 这给出图
\[
\begin{tikzcd}[column sep=small]
x_0\arrow[r]\arrow[d]&x\arrow[d]\arrow[dr]&\\
\tau\arrow[r]\arrow[rr,bend right=20]&y\arrow[r,dashed,"\beta"]&y'
\end{tikzcd}
\qquad\text{在}\qquad
\begin{tikzcd}[column sep=small]
\Spec\Bbbk\arrow[r,hook]\arrow[d]&\Spec A\arrow[d]\\
\Spec\Bbbk[I]\arrow[r,hook]&\Spec(\Bbbk[I]\times_{\Bbbk}A).
\end{tikzcd}
\]
由 $(\mathrm{RS}_2)$ 的第二部分, 存在补全该图的 $\beta:y\to y'$, 故 $\mathcal X(A)$ 中的 $p_1^*y$ 和 $p_1^*y'$ 定义 $\operatorname{Lift}_{x_0}(A)$ 中的同一元素. 最后, 若 $(\mathrm{RS}_1)$ 成立 (相应地, $(\mathrm{RS}_4)$ 成立), 则 $\mathcal X(A\times_{A_0}A)\to\mathcal X(A)\times_{\mathcal X(A_0)}\mathcal X(A)$ 本质满射 (相应地, 为等价), 因而作用传递 (相应地, 自由且传递).
\end{proof}
\begin{proof}[定理 \bookref{C.4.6} 的证明]
我们建立判据的充分性, 必要性留给读者. 切空间 $T_{\mathcal X}:=\mathcal X(\Bbbk[\epsilon])/{\sim}$ 由引理 \bookref{C.4.9}(1) 具有向量空间结构, 并由 $(\mathrm{RS}_2)$ 有限维. 令 $N=\dim_{\Bbbk}T_{\mathcal X}$, 取基 $x_1,\ldots,x_N\in T_{\mathcal X}$. 定义 $S=\Bbbk[\![x_1,\ldots,x_N]\!]$, 极大理想为 $\mathfrak m_S=(x_1,\ldots,x_n)$. 我们归纳地构造下降理想列 $J_0\supseteq J_1\supseteq\cdots$ 及对象 $\eta_n\in\mathcal X(S/J_n)$, 连同 $\Spec S/J_n\hookrightarrow\Spec S/J_{n+1}$ 上的映射 $\eta_n\to\eta_{n+1}$. 令 $J_0=\mathfrak m_S$, $\eta_0=x_0\in\mathcal X(\Bbbk)$. 又令 $J_1=\mathfrak m_S^2$, 故 $S/J_1\simeq\Bbbk[T_{\mathcal X}]$. 由引理 \bookref{C.4.9}(1) 的双射 $F_{\mathcal X}(\Bbbk[T_{\mathcal X}])\simeq T_{\mathcal X}\otimes_{\Bbbk}T_{\mathcal X}$, 元素 $\sum_i x_i\otimes x_i$ 定义对象 $\eta_1\in\mathcal X(S/J_1)$ 的同构类, 使得诱导的映射 $\Spec S/J_1\to\mathcal X$ 在切空间上诱导双射. 按构造, 在 $\Spec\Bbbk\hookrightarrow\Spec S/J_1$ 上有映射 $\eta_0\to\eta_1$.

假设已构造 $J_n$ 及 $\eta_{n-1}\to\eta_n$. 断言理想集
\begin{equation}\label{src:C.4.11}
\Sigma=\left\{J\subseteq S\ \middle|\ \begin{gathered}
\mathfrak m_SJ_n\subseteq J\subseteq J_n,\ \text{且存在 }\eta_n\to\eta\\
\text{位于 }\Spec S/J_n\hookrightarrow\Spec S/J\text{ 上}
\end{gathered}\right\}
\end{equation}
具有极小元素. 事实上, 因为 $J_n\in\Sigma$, 此集合非空, 所以只需检验对 $J,K\in\Sigma$ 有 $J\cap K\in\Sigma$. 注意 $J_n/\mathfrak m_SJ_n$ 为 $\Bbbk$-向量空间, 子空间为 $J/\mathfrak m_SJ_n$ 和 $K/\mathfrak m_SJ_n$. 因而可选取理想 $J\subseteq J'\subseteq J_n$, 满足 $J\cap K=J'\cap K$ 且 $J'+K=J_n$. 有图
\[
\begin{tikzcd}
\eta_{J\cap K}&\eta_J|_{S/J'}\arrow[l,dashed]\\
\eta_K\arrow[u,dashed]&\eta_n\arrow[l]\arrow[u]
\end{tikzcd}
\qquad\text{在}\qquad
\begin{tikzcd}
S/(J\cap K)\arrow[r]\arrow[d]&S/J'\arrow[d]\\
S/K\arrow[r]\arrow[ur,phantom,"\square",very near start]&S/J_n,
\end{tikzcd}
\]
其中 $\eta_J\in\mathcal X(S/J)$ 和 $\eta_K\in\mathcal X(S/K)$ 为对应 $J$ 和 $K$ 的对象. 条件 $(\mathrm{RS}_1)$ 蕴含可以补全该图, 因而 $J\cap K\in\Sigma$.

定义 $J=\bigcap_nJ_n$, $R=S/J$, $I_n=J_n/J$. 断言对每个 $n$ 存在 $N_n$, 使 $J_{N_n}\subseteq\mathfrak m_S^{n+1}+J$; 换言之, $(I_n)$ 在 $R$ 上定义的拓扑为 $\mathfrak m_R$-进拓扑. 对每个 $n$, 因 $S/(\mathfrak m_S^n)$ 为 Artin 环, 存在 $N_n$, 满足 $J_{N_n}+\mathfrak m_S^n=J_{N_{n+1}}+\mathfrak m_S^n=\cdots$. 令
\[
E=\lim_n(J_{N_n}+\mathfrak m_S^n)/\mathfrak m_S^n\subseteq\lim_nS/\mathfrak m_S^n=S.
\]
断言 $E\subseteq J$: 对 $f\in E$ 及任意 $m$, 当 $n\gg0$ 时有 $f\in J_m+\mathfrak m_S^n$, 由 Krull 交定理推出断言. 因为 $E$ 满射到 $(J_{N_n}+\mathfrak m_S^{n+1})/(\mathfrak m_S^{n+1})$, 自然映射 $J\to(J_{N_n}+\mathfrak m_S^{n+1})/(\mathfrak m_S^{n+1})$ 也满射, 断言得证. 由 $I_{N_n}\subseteq\mathfrak m_R^{n+1}$, 可定义 $\xi_n:=\eta_{N_n}|_{R/\mathfrak m_R^{n+1}}$.

还需证明 $R$ 上的形式形变 $\xi:=\{\xi_n\}$ 万有. 假设 $B\twoheadrightarrow A$ 为 $\mathrm{Art}_{\Bbbk}$ 中核 $\Bbbk$ 的满射, 且有图
\[
\begin{tikzcd}
x\arrow[r]\arrow[d]&\xi\\
y\arrow[ur,dashed]&
\end{tikzcd}
\qquad\text{在}\qquad
\begin{tikzcd}
\Spec A\arrow[r,"g"]\arrow[d,hook]&h_R\\
\Spec B\arrow[ur,dashed,"\widetilde g"']&
\end{tikzcd}
\]
必须构造延拓 $x\to\xi$ 的态射 $y\to\xi$. 断言只需构造延拓 $g$ 的态射 $\widetilde g:\Spec B\to h_R$ (即环映射 $R\to B$). 因为 $h_R(\Bbbk[\epsilon])\to T_{\mathcal X}$ 双射, 引理 \bookref{C.4.9}(2) 给出 $T_{\mathcal X}$ 在 $\operatorname{Lift}_x(B)$ 和 $\operatorname{Lift}_g(B)$ 上的作用; 它们分别为 $x$ 和 $g$ 提升到 $\mathcal X(B)$ 和 $h_R(B)$ 中对象的同构类集合. 又因为 $\mathcal X$ 满足 $(\mathrm{RS}_1)$, 在 $\operatorname{Lift}_x(B)$ 上的作用传递. 因而可找到 $\tau\in T_{\mathcal X}$, 满足 $y=\tau\cdot(\widetilde g^*\xi)=(\tau\cdot\widetilde g)^*\xi$. 这给出 $\tau\cdot\widetilde g:\Spec B\to h_R$ 上的箭头 $y\to\xi$.

为构造 $\widetilde g$, 选取 $n$, 使 $R\to A$ 分解为 $R\to R/I_n=S/J_n\to A$. 只需证明 $\Spec A\to\Spec S/J_n$ 延拓为映射 $\Spec B\to\Spec S/J_{n+1}$; 为此, 只需存在使图
\[
\begin{tikzcd}[column sep=small]
\Spec A\arrow[r]\arrow[d,hook]&\Spec S/J_n\arrow[d,hook]\arrow[dr,hook]&\\
\Spec B\arrow[r]&\Spec(B\times_A(S/J_n))\arrow[r,dashed]&\Spec S/J_{n+1}
\end{tikzcd}
\]
交换的虚线箭头. 注意由 $\ker(B\to A)=\Bbbk$, 有 $\ker(B\times_A(S/J_n)\to S/J_n)=\Bbbk$. 因为 $S$ 为幂级数环, 可选取 $S\to S/J_n\to A$ 的延拓 $S\to B$. 它诱导映射 $S\to B\times_A(S/J_n)$. 若此映射不满射, 则 $J_n$ 的每个元素必映到 $B\times_A(S/J_n)$ 中的 $0$, 故 $B\times_A(S/J_n)\to S/J_n$ 具有截面, 给出所需提升. 否则, $B\times_A(S/J_n)=S/K$, 其中 $K=\ker(S\to B\times_A(S/J_n))$. 理想 $K$ 属于 \bookeqref{C.4.11} 所定义的理想集: 包含关系 $K\subset J_n$ 显然, $\mathfrak m_SJ_n\subseteq K$ 由等式 $\ker(B\to A)=\Bbbk$ 推出, 而 $\Spec S/J_n\hookrightarrow\Spec S/K$ 上的 $\eta_n\to\eta$ 的存在由对上述方块应用 $(\mathrm{RS}_1)$ 推出. 由 $J_{n+1}$ 的极小性, 有 $J_{n+1}\subseteq K$, 故有环映射 $S/J_{n+1}\to S/K=B\times_A(S/J_n)$, 诱导所需虚线箭头.

最后, 假设 $\mathcal X$ 等价于函子 $F$, 且 $(\mathrm{RS}_4)$ 成立. 给定核 $\Bbbk$ 的满射 $B\twoheadrightarrow A$ 及 $x\in F(A)$, 需证明每个图
\[
\begin{tikzcd}
\Spec A\arrow[r,"g"]\arrow[d,hook]&h_R\arrow[d]\\
\Spec B\arrow[r]\arrow[ur,dashed]&F
\end{tikzcd}
\]
中存在唯一提升. 由引理 \bookref{C.4.9}(2), 映射 $\operatorname{Lift}_g(B)\to\operatorname{Lift}_x(B)$ 双射, 因为二者均为 $T_{\mathcal X}$ 下的 torsor. 因而存在唯一提升. 另见 \cite[定理 $2.11$]{Sch68}, \cite[定理 VI.$1.11$]{SGA7-I}, \cite[定理 $16.2$]{Har10}, \cite[定理 $2.3.2$]{Ser06} 及 \cite[Tag 06IX]{SP}.
\end{proof}
